\documentclass[11pt]{article}
\usepackage[T1]{fontenc}
\usepackage[utf8]{inputenc}
\usepackage{lmodern}
\usepackage[margin=1in]{geometry}
\usepackage{amsmath,amssymb,amsthm,mathtools,bm}
\usepackage{longtable,booktabs,array,multirow}
\usepackage{textcomp}
\usepackage{graphicx}
\usepackage{enumitem}
\usepackage{xcolor}
\usepackage{framed}
\usepackage{fancyvrb}
\usepackage{ragged2e}
\usepackage{hyperref}
\usepackage{xurl}
\newcommand{\OPHPaperReleaseID}{r1412}
\newcommand{\OPHPaperReleaseDate}{May 4, 2026}
\newcommand{\OPHPaperReleaseBanner}{%
\begin{center}
\small\textbf{Paper release:} \texttt{\OPHPaperReleaseID}\quad\textbf{Released:} \OPHPaperReleaseDate
\end{center}
}


% Springer/pdflatex compatibility: provide fallbacks for Unicode math/text glyphs.
\newcounter{none}
\providecommand{\theHnone}{\arabic{none}}
\providecommand{\textsubscript}[1]{$_{\mathrm{#1}}$}
\newcommand{\ophid}[1]{{\small\path{#1}}}
\ifdefined\DeclareUnicodeCharacter
\DeclareUnicodeCharacter{00B1}{\ensuremath{\pm}}
\DeclareUnicodeCharacter{00B2}{\textsuperscript{2}}
\DeclareUnicodeCharacter{00B3}{\textsuperscript{3}}
\DeclareUnicodeCharacter{00B7}{\ensuremath{\cdot}}
\DeclareUnicodeCharacter{00B9}{\textsuperscript{1}}
\DeclareUnicodeCharacter{00D7}{\ensuremath{\times}}
\DeclareUnicodeCharacter{02B8}{\textsuperscript{y}}
\DeclareUnicodeCharacter{0393}{\ensuremath{\Gamma}}
\DeclareUnicodeCharacter{0394}{\ensuremath{\Delta}}
\DeclareUnicodeCharacter{039B}{\ensuremath{\Lambda}}
\DeclareUnicodeCharacter{03A3}{\ensuremath{\Sigma}}
\DeclareUnicodeCharacter{03A6}{\ensuremath{\Phi}}
\DeclareUnicodeCharacter{03A9}{\ensuremath{\Omega}}
\DeclareUnicodeCharacter{03B1}{\ensuremath{\alpha}}
\DeclareUnicodeCharacter{03B2}{\ensuremath{\beta}}
\DeclareUnicodeCharacter{03B3}{\ensuremath{\gamma}}
\DeclareUnicodeCharacter{03B4}{\ensuremath{\delta}}
\DeclareUnicodeCharacter{03B5}{\ensuremath{\varepsilon}}
\DeclareUnicodeCharacter{03B7}{\ensuremath{\eta}}
\DeclareUnicodeCharacter{03B8}{\ensuremath{\theta}}
\DeclareUnicodeCharacter{03BA}{\ensuremath{\kappa}}
\DeclareUnicodeCharacter{03BB}{\ensuremath{\lambda}}
\DeclareUnicodeCharacter{03BC}{\ensuremath{\mu}}
\DeclareUnicodeCharacter{03BD}{\ensuremath{\nu}}
\DeclareUnicodeCharacter{03BE}{\ensuremath{\xi}}
\DeclareUnicodeCharacter{03C0}{\ensuremath{\pi}}
\DeclareUnicodeCharacter{03C1}{\ensuremath{\rho}}
\DeclareUnicodeCharacter{03C3}{\ensuremath{\sigma}}
\DeclareUnicodeCharacter{03C4}{\ensuremath{\tau}}
\DeclareUnicodeCharacter{03C6}{\ensuremath{\phi}}
\DeclareUnicodeCharacter{03C7}{\ensuremath{\chi}}
\DeclareUnicodeCharacter{03C8}{\ensuremath{\psi}}
\DeclareUnicodeCharacter{03C9}{\ensuremath{\omega}}
\DeclareUnicodeCharacter{1D9C}{\textsuperscript{c}}
\DeclareUnicodeCharacter{2020}{\ensuremath{\dagger}}
\DeclareUnicodeCharacter{2070}{\textsuperscript{0}}
\DeclareUnicodeCharacter{2074}{\textsuperscript{4}}
\DeclareUnicodeCharacter{2075}{\textsuperscript{5}}
\DeclareUnicodeCharacter{2076}{\textsuperscript{6}}
\DeclareUnicodeCharacter{2077}{\textsuperscript{7}}
\DeclareUnicodeCharacter{2078}{\textsuperscript{8}}
\DeclareUnicodeCharacter{2079}{\textsuperscript{9}}
\DeclareUnicodeCharacter{207A}{\textsuperscript{+}}
\DeclareUnicodeCharacter{207B}{\textsuperscript{-}}
\DeclareUnicodeCharacter{207D}{\textsuperscript{(}}
\DeclareUnicodeCharacter{207E}{\textsuperscript{)}}
\DeclareUnicodeCharacter{207F}{\textsuperscript{n}}
\DeclareUnicodeCharacter{2080}{\textsubscript{0}}
\DeclareUnicodeCharacter{2081}{\textsubscript{1}}
\DeclareUnicodeCharacter{2082}{\textsubscript{2}}
\DeclareUnicodeCharacter{2083}{\textsubscript{3}}
\DeclareUnicodeCharacter{2085}{\textsubscript{5}}
\DeclareUnicodeCharacter{2086}{\textsubscript{6}}
\DeclareUnicodeCharacter{2087}{\textsubscript{7}}
\DeclareUnicodeCharacter{2088}{\textsubscript{8}}
\DeclareUnicodeCharacter{208B}{\textsubscript{-}}
\DeclareUnicodeCharacter{2102}{\ensuremath{\mathbb{C}}}
\DeclareUnicodeCharacter{210B}{\ensuremath{\mathcal{H}}}
\DeclareUnicodeCharacter{2113}{\ensuremath{\ell}}
\DeclareUnicodeCharacter{2124}{\ensuremath{\mathbb{Z}}}
\DeclareUnicodeCharacter{2192}{\ensuremath{\to}}
\DeclareUnicodeCharacter{2193}{\ensuremath{\downarrow}}
\DeclareUnicodeCharacter{2194}{\ensuremath{\leftrightarrow}}
\DeclareUnicodeCharacter{21D2}{\ensuremath{\Rightarrow}}
\DeclareUnicodeCharacter{2202}{\ensuremath{\partial}}
\DeclareUnicodeCharacter{2208}{\ensuremath{\in}}
\DeclareUnicodeCharacter{220F}{\ensuremath{\prod}}
\DeclareUnicodeCharacter{2212}{\ensuremath{-}}
\DeclareUnicodeCharacter{221A}{\ensuremath{\surd}}
\DeclareUnicodeCharacter{221D}{\ensuremath{\propto}}
\DeclareUnicodeCharacter{221E}{\ensuremath{\infty}}
\DeclareUnicodeCharacter{2229}{\ensuremath{\cap}}
\DeclareUnicodeCharacter{222B}{\ensuremath{\int}}
\DeclareUnicodeCharacter{223C}{\ensuremath{\sim}}
\DeclareUnicodeCharacter{2245}{\ensuremath{\cong}}
\DeclareUnicodeCharacter{2248}{\ensuremath{\approx}}
\DeclareUnicodeCharacter{2261}{\ensuremath{\equiv}}
\DeclareUnicodeCharacter{2264}{\ensuremath{\leq}}
\DeclareUnicodeCharacter{2265}{\ensuremath{\geq}}
\DeclareUnicodeCharacter{226A}{\ensuremath{\ll}}
\DeclareUnicodeCharacter{2272}{\ensuremath{\lesssim}}
\DeclareUnicodeCharacter{2295}{\ensuremath{\oplus}}
\DeclareUnicodeCharacter{2297}{\ensuremath{\otimes}}
\DeclareUnicodeCharacter{2609}{\ensuremath{\odot}}
\DeclareUnicodeCharacter{2713}{\ensuremath{\checkmark}}
\DeclareUnicodeCharacter{27E8}{\ensuremath{\langle}}
\DeclareUnicodeCharacter{27E9}{\ensuremath{\rangle}}
\DeclareUnicodeCharacter{27F9}{\ensuremath{\Longrightarrow}}
\fi

\hypersetup{colorlinks=true,linkcolor=blue,citecolor=blue,urlcolor=blue}
\setcounter{tocdepth}{2}
\setcounter{secnumdepth}{3}
\setlength{\LTleft}{0pt}
\setlength{\LTright}{0pt}
\setlength{\tabcolsep}{1.2pt}
\setlength{\emergencystretch}{3em}
\renewcommand{\arraystretch}{1.05}

\providecommand{\tightlist}{%
 \setlength{\itemsep}{0pt}\setlength{\parskip}{0pt}}
\definecolor{shadecolor}{RGB}{248,248,248}
\newenvironment{Shaded}{\begin{snugshade}}{\end{snugshade}}
\DefineVerbatimEnvironment{Highlighting}{Verbatim}{commandchars=\\\{\}}
\newcommand{\NormalTok}[1]{#1}
\newcommand{\BuiltInTok}[1]{#1}
\newcommand{\CommentTok}[1]{#1}
\newcommand{\ExtensionTok}[1]{#1}
\newcommand{\AttributeTok}[1]{#1}
\newcommand{\SU}{\mathrm{SU}}
\newcommand{\U}{\mathrm{U}}

\title{Observers Are All You Need}
\author{B. M\"uller \and Peter Nguyen}
\date{\OPHPaperReleaseDate}

\begin{document}
\maketitle
\OPHPaperReleaseBanner

\begin{abstract}
Observer-Patch Holography (OPH) asks whether familiar physics can be rebuilt from a very small
starting point: no observer sees the whole world at once, each observer has access only to local
data, and neighboring descriptions must agree where they overlap. In that picture, spacetime,
gauge structure, particles, and classical records are not put in by hand at the start. They are
features of whatever survives global consistency among many partial viewpoints on a finite-capacity
screen.

This paper is the synthesis surface for the OPH suite. It pulls together the fixed-cutoff
collar and higher-gauge theorem packages, the consensus/fixed-point formulation of overlap repair,
the conditional route to Lorentzian geometry and Einstein dynamics, the compact-gauge route to the
realized Standard Model quotient
\(\SU(3)\times\SU(2)\times\U(1)/\mathbb Z_6\), the particle-spectrum lane, and the screen
microphysics and observer lane. It presents, in one place, the theorem-bearing results, the
conditional branches, and the explicit external objects above the recovered core.

At fixed cutoff, OPH has a sharp local picture: patches glue, higher-gauge defects are
controlled, physical UV uniqueness is quotient-level and ancilla-stable, and the microphysics lane
has an explicit regulated architecture together with closed measurement and
checkpoint/restoration packages. On the continuum side, OPH gives a conditional
Lorentzian/Jacobson-type gravity branch and a conditional compact-gauge route whose realized
low-energy output is the Standard Model quotient above with the exact hypercharge lattice and the
counting chain \(N_g=3\) then \(N_c=3\). On the particle side, the structural massless carriers
are fixed, the electroweak lane is closed at theorem level, the Higgs/top stage is emitted on its
secondary branch, and the neutrino lane matches the observed mixing/hierarchy pattern while
leaving one reduced normalization invariant open; the charged-lepton, quark, and hadron lanes keep
their remaining closure objects explicit. The main external burden is the continuum/BW lift on the
geometric subnet: extract the prime geometric cap pair from transported overlap data, then prove
ordered cut-pair rigidity on that extracted limit.
\end{abstract}

\tableofcontents
\newpage

\section{Introduction}

This synthesis paper reads the OPH suite on one unified synthesis surface. It brings
together the recovered core from Ref.~\cite{ophcompact}, the finite patch-net consensus spine from
Ref.~\cite{ophconsensus}, the screen-microphysics engineering lane from Ref.~\cite{ophmicrophysics},
and the downstream particle-spectrum lane from Ref.~\cite{ophparticles}. The guiding
question is whether overlap consistency for local observer descriptions on a finite-capacity screen
is strong enough to recover the structural features of low-energy physics, and if so, how much of
that recovery is theorem-level and how much depends on explicit scaling-limit,
categorical, or continuation premises.

The answer is sharply tiered. Fixed-cutoff collar structure, higher-gauge defect transport,
the consensus/fixed-point formulation, the quotient-level UV uniqueness statement, the Born-rule
record package, and the checkpoint/restoration package are theorem-bearing on their stated finite
surfaces. The Lorentz/null-modular/Einstein chain and compact gauge reconstruction remain
conditional continuum or categorical branches. The particle lane is the most tiered part of
the suite: the structural carrier skeleton and realized Standard Model branch are fixed, the
electroweak stage is closed, the Higgs/top stage is a secondary quantitative branch, the neutrino
lane is shape-accurate with one reduced normalization invariant open, and the
charged-lepton, quark, and hadron lanes remain explicitly unfinished for different reasons.

The paper is organized by lane. Section~2 gives the synthesis-level map. Section~3 carries the
recovered-core trunk: axioms, premise ledger, collar tools, gluing, the Lorentz/gravity branch,
and the compact-gauge/Standard-Model branch. Section~4 records the consensus lane. Section~5
summarizes the particle lane. Section~6 records the screen-microphysics and observer lane.
Section~7 records the conditional worldsheet lane. Section~8 gives the cross-lane theorem
boundary and open directions. The appendices collect extra architecture and interpretive
material that is useful in the synthesis paper but not part of the recovered core itself.

\section{Overview of Results}

The OPH suite supports the following lane-level split.

\begin{enumerate}[leftmargin=2em]
\item \textbf{At fixed cutoff, local structure and repair are sharp.} Technically, the
collar theorem package is explicit and finite-range, the genuinely noncentral higher-gauge branch
is closed at fixed cutoff, and the consensus lane gives a unique schedule-independent normal form
on the gauge quotient under the stated finite patch-net hypotheses.
\item \textbf{Relativity appears as a controlled continuum branch.}
Technically, the Lorentz/null-modular/Einstein chain is a refinement/scaling-limit branch with
explicit canonical technical premises.
\item \textbf{The Standard Model gauge skeleton appears as a constrained consistency output.}
Technically, compact gauge reconstruction and the realized
\(\SU(3)\times\SU(2)\times\U(1)/\mathbb Z_6\) branch follow under the stated transportability and
categorical premises, with the realized counting chain \(N_g=3\) then \(N_c=3\).
\item \textbf{The particle lane has real outputs and a nonuniform closure profile.}
Technically, the structural carriers are exact, the electroweak stage is closed, the Higgs/top
stage is quantitatively emitted on a secondary branch, and the charged-lepton, quark, neutrino,
and hadron lanes carry named remaining closure objects.
\item \textbf{The screen-microphysics lane is explicit and concrete.} Technically, it
contains an explicit regulated reference architecture, sharpened patch-net and edge-sector
branches, a closed fixed-cutoff Born-rule package, and a closed checkpoint/restoration package.
\item \textbf{The UV picture is sharper at fixed cutoff than in the continuum.} Technically, OPH
fixes the physical UV branch only modulo implementation hiding and inert ancillary refinement; the
UV burden is the continuum BW/geometric lift and the refinement-limit transportable-sector
lift.
\item \textbf{The string/worldsheet lane remains conditional.} Technically, it is carried only as
a continuation of the OPH heat-kernel edge-sector theorem, not as part of the recovered core.
\end{enumerate}

\subsection*{Companion Papers}

The detailed depth surfaces used throughout this synthesis are:
\begin{enumerate}[leftmargin=2em]
\item \emph{A Conditional Reconstruction Program for Relativity and Standard Model Structure from Observer-Overlap Consistency: Observers Are All You Need (Compact)}~\cite{ophcompact}, which carries the premise ledger, the recovered-core Lorentz/gravity route, and the compact-gauge/Standard-Model branch.
\item \emph{Reality as a Consensus Protocol: The Fixed-Point Computation That Implements Physics}~\cite{ophconsensus}, which carries the finite patch-net fixed-point, defect, quotient, and record-layer consensus package.
\item \emph{Screen Microphysics, Patches, and Observer Synchronization in OPH}~\cite{ophmicrophysics}, which carries the regulated reference architecture together with the measurement and observer checkpoint/restoration packages.
\item \emph{Deriving the Particle Zoo from Observer Consistency}~\cite{ophparticles}, which carries the structural-to-family particle route and the per-sector claim-tier ledger.
\end{enumerate}

\def\OPHSkipEmbeddedReferences{1}
\input{tex_fragments/PAPER.tex}
\section{Consensus, Defects, and Implementation Hiding}

The OPH core admits an equivalent finite patch-net formulation. Local patch descriptions agree on
overlaps, local repair rules remove mismatches, and the physical output is the
schedule-independent normal form on the gauge quotient. This section records the fixed-point,
defect, and gauge-quotient statements in that language.

The consensus lane also carries a simple law-selection model: once one equips candidate repair
rules with a fitness functional, replicator dynamics gives a clean toy picture in which more
successful reconciliation laws dominate a competing pool. That selection picture is useful for the
interpretation of OPH, but it is not part of the recovered-core gravity/gauge theorem surface.

\paragraph{Theorem (Asynchronous confluence).}
On a finite patch net with local repair maps, if every accepted repair strictly decreases a
finite inconsistency potential and the repair relation is locally confluent, then every initial
configuration has a unique schedule-independent normal form.

\paragraph{Proof sketch.}
Strict decrease of the inconsistency potential rules out infinite repair sequences. Newman's
lemma~\cite{newman42} then upgrades local confluence to global confluence. Therefore every initial state has a
unique normal form, independent of the asynchronous update schedule.

\paragraph{Theorem (Cycle holonomy and higher-gauge defects).}
Pairwise overlap agreement is not enough for global consistency. On the abelian branch, a global
solution exists iff the cycle holonomy vanishes on every loop. On the genuinely noncentral
branch, weak gluing data define a crossed-module Cech class
\[
q_\Sigma \in \check H^2(N_\Sigma,H_\Sigma \to G_\Sigma),
\]
and strict global reconciliation exists iff \(q_\Sigma=0\). Nonzero \(q_\Sigma\) labels stable
fixed-cutoff higher-gauge defects.

\paragraph{Proof sketch.}
In the abelian case, summing edge constraints around a loop telescopes, so vanishing of the loop
sum is necessary and sufficient for path-independent reconstruction. In the crossed-module case,
local rechartings are exactly crossed-module coboundaries, so they preserve the class and remove
the defect precisely when the class is trivial.

\paragraph{Theorem (Gauge quotient, ancilla-stable physical uniqueness, and record closure).}
If the repair dynamics is gauge covariant, the normal-form map descends to the quotient by
implementation hiding:
\[
[s]\mapsto [\operatorname{nf}(s)].
\]
Gauge-invariant observables are therefore uniquely fixed on the quotient, and inert ancillary
refinement does not change the physical output. On the classical record layer, commuting closure
operators on a finite lattice converge under any asynchronous fair schedule to the least common
fixed point above the initial record state.

\paragraph{Proof sketch.}
Gauge covariance carries each repair sequence to a repair sequence on a gauge-related state, so
uniqueness of normal forms forces equivariance of the normal-form map. Inert ancillary factors
remain spectators under the lifted dynamics, so gauge-invariant observables are unchanged. For
records, commutativity and idempotence make the closure order-independent and force convergence to
the least common fixed point.

\paragraph{Extended derivation.}
The full consensus, gauge-quotient, and record-layer theorem package is developed in
\emph{Reality as a Consensus Protocol: The Fixed-Point Computation That Implements Physics}~\cite{ophconsensus}.

\section{Particle-Spectrum Branch}

The particle branch follows the same logical order as
\emph{Deriving the Particle Zoo from Observer Consistency}~\cite{ophparticles}.
One first fixes the OPH basis and the compact gauge branch, then identifies the realized Standard
Model quotient, and only then asks how transport-stable overlap data organize the observed
particle families.

\paragraph{Structural particle theorem package.}
On the realized OPH branch, the structural carrier skeleton is fixed by the compact gauge and
MAR-admissibility chain:
\[
\frac{\SU(3)\times\SU(2)\times\U(1)}{\mathbb Z_6},
\qquad
N_g=3,
\qquad
N_c=3.
\]
This determines the exact structural zero rows carried by the gauge and gravity sectors:
the photon, gluons, and graviton are massless structural carriers on the realized branch.

\paragraph{Proof route.}
The proof route is the one developed in
\emph{A Conditional Reconstruction Program for Relativity and Standard Model Structure from Observer-Overlap Consistency: Observers Are All You Need (Compact)}~\cite{ophcompact}
and
\emph{Deriving the Particle Zoo from Observer Consistency}~\cite{ophparticles}.
Refinement-stable transportable edge sectors reconstruct a compact internal symmetry group under
the stated categorical premises. MAR then selects the realized admissible branch, anomaly freedom
fixes the hypercharge lattice, and the realized one-generation/one-Higgs package forces \(N_g=3\)
and then \(N_c=3\). The particle branch builds on that structural skeleton and does not bypass
it.

\scriptsize
\setlength{\tabcolsep}{2pt}
\begin{longtable}{@{}>{\raggedright\arraybackslash}p{0.175\textwidth}>{\raggedright\arraybackslash}p{0.165\textwidth}>{\raggedright\arraybackslash}p{0.225\textwidth}>{\raggedright\arraybackslash}p{0.325\textwidth}@{}}
\toprule
Sector & Claim tier & What is fixed here & Boundary to stronger claim \\
\midrule
\endfirsthead
\toprule
Sector & Claim tier & What is fixed here & Boundary to stronger claim \\
\midrule
\endhead
\bottomrule
\endfoot
Structural carriers & realized-branch structural theorem & exact \(m_\gamma=m_g=m_{\mathrm{grav}}=0\); realized quotient; \(N_g=3\); \(N_c=3\) & theorem-grade structural exactness on the particle side \\
Electroweak bosons & calibration branch + exact sidecar & exact frozen-repair pair \(M_W=80.377~\mathrm{GeV}\), \(M_Z=91.18797809193725~\mathrm{GeV}\) & exact only on the frozen authoritative repair surface; the public theorem rows remain the target-free electroweak outputs \\
Higgs/top stage & secondary quantitative branch + closed forward seed + exact sidecars & closed one-scalar forward seed for the public D11 rows together with an exact Higgs/top inverse pair & the public D11 rows live on a secondary-quantitative branch; the inverse pair is compare-only and does not replace the forward rows \\
Quark family & continuation branch + exact witness & exact same-family quark sextet on the same ordered three-point family, with the ordered three-point quadratic readout closed on that scope & the selected quark sheet is the wrong branch; the broader honest frontier is the light-quark overlap-defect scalar \(\Delta_{ud}^{\mathrm{overlap}}\), and on the emitted D12 mass ray this is equivalently \ophid{quark_d12_t1_value_law}, with \ophid{intrinsic_scale_law_D12} as the derived wrapper; a continuation-only D12 internal backread sidecar fixes the mass-side scalar package numerically without removing the wrong CKM branch \\
Charged leptons & continuation branch + exact witness & exact same-family charged triple plus exact same-carrier centered readback, with the ordered three-point quadratic readout closed on that scope & same-family-only on the exact witness; theorem lane waits on \(\widehat C_e^{\mathrm{cand}}\) and the post-promotion lift whose descended scalar is \(\mu_{\mathrm{phys}}(Y_e)\), with the physical identity-mode equalizer beneath that scalar \\
Neutrinos & weighted-cycle continuation + exact adapter & exact compare-only two-parameter masses/splittings on the explicit positive selector segment & one reduced bridge-correction invariant \(C_\nu\) remains open above the emitted proxy, so the exact adapter stays compare-only \\
Hadrons & nonperturbative continuation & stable-channel and readout architecture are defined & one production backend export bundle, then executed unquenching, runtime receipt, and production systematics \\
\end{longtable}
\normalsize

\paragraph{Claim-tier boundary.}
Structural carrier content and the realized Standard Model branch are theorem-bearing outputs.
Electroweak calibration, the Higgs/top stage, quark rows, charged-lepton rows, neutrino rows, and
hadron rows sit on their stated calibration, continuation, or nonperturbative claim tiers.

\paragraph{Exact non-hadron hit surface.}
For quick orientation, the exact non-hadron lanes are:

\begin{table}[h]
\centering
\scriptsize
\setlength{\tabcolsep}{3pt}
\begin{tabular}{@{}>{\raggedright\arraybackslash}p{0.14\textwidth}>{\raggedright\arraybackslash}p{0.17\textwidth}>{\raggedright\arraybackslash}p{0.31\textwidth}>{\raggedright\arraybackslash}p{0.30\textwidth}@{}}
\toprule
Lane & Exact output(s) & Exact OPH chain on the paper surface & Caveat \\
\midrule
Structural carriers & \(m_\gamma=m_g=m_{\mathrm{grav}}=0\) & axioms \(\rightarrow\) realized electromagnetic/color/dynamical-metric carrier skeleton \(\rightarrow\) symmetry-protected zeros & theorem-grade structural exactness \\
Electroweak exact sidecar & exact frozen \(W/Z\) pair & axioms \(\rightarrow D7\rightarrow D10\) calibration chain \(\rightarrow\) exact frozen repair pair & exact only on the frozen authoritative repair surface; compare-only beneath the target-free theorem \\
Higgs/top exact sidecar & exact Higgs/top inverse pair & axioms \(\rightarrow D10\) gauge core \(\rightarrow D11\) Jacobian \(\rightarrow\) exact inverse slice & compare-only inverse slice; the public D11 rows come from the closed one-scalar forward seed \\
Charged exact witness & exact same-family charged triple & axioms \(\rightarrow\) shared excitation dictionary \(\rightarrow\) ordered charged carrier \(\rightarrow\) exact centered readback \(\rightarrow\) closed quadratic readout theorem \(\rightarrow\) same-family exact witness & same-family-only; theorem lane waits on \(\widehat C_e^{\mathrm{cand}}\) and the post-promotion lift whose descended scalar is \(\mu_{\mathrm{phys}}(Y_e)\), with the physical identity-mode equalizer beneath that scalar \\
Quark exact witness & exact same-family quark sextet & axioms \(\rightarrow\) shared excitation dictionary \(\rightarrow\) D12 mass-side continuation \(\rightarrow\) selected D12 sheet \(\rightarrow\) closed quadratic readout theorem \(\rightarrow\) same-family exact witness & same-family-only; the selected D12 sheet is the wrong CKM branch, the broader honest frontier is the light-quark overlap-defect scalar \(\Delta_{ud}^{\mathrm{overlap}}\), and on the emitted D12 mass ray this is equivalently \ophid{quark_d12_t1_value_law}, with \ophid{intrinsic_scale_law_D12} as the derived wrapper; a continuation-only D12 internal backread sidecar fixes the mass-side scalar package numerically without removing the wrong CKM branch \\
Neutrino exact adapter & exact compare-only mass triple and representative central splittings & axioms \(\rightarrow\) same-label scalar certificate \(\rightarrow\) weighted-cycle branch \(\rightarrow\) explicit positive selector segment \(\rightarrow\) two-parameter exact adapter & compare-only; the theorem lane waits on the reduced bridge-correction invariant \(C_\nu\) \\
\bottomrule
\end{tabular}
\end{table}

\noindent The explicit values behind these exact-hit lanes are the exact frozen
\(W/Z\) pair, the exact Higgs/top inverse pair, the exact same-family charged triple,
the exact same-family quark sextet, and the exact compare-only neutrino mass triple
with exact representative central splittings recorded in Ref.~\cite{ophparticles}'s exact-output surface. The public Higgs/top rows come from the closed one-scalar D11 forward seed rather than from the compare-only inverse slice.

\paragraph{Extended derivation.}
The structural-to-family route and the full per-sector claim-tier audit are developed in
\emph{Deriving the Particle Zoo from Observer Consistency}~\cite{ophparticles}.

\section{Screen Microphysics, Records, and Observer Continuation}

The screen-microphysics lane is the engineering side of the OPH suite. It asks
whether some microscopic model might exist in principle and gives one explicit regulated
reference architecture: a finite gauge-register or quantum-link style screen model in which local
registers, patches, overlaps, edge observables, records, repair instruments, and observer-facing
interfaces are all made concrete at fixed cutoff.

This lane has four theorem-bearing components. First, the reference architecture itself is
explicit. Second, the regulated patch-net embedding branch and the edge-sector thermal/Casimir
branch are both sharpened, each with one named closure pass above the declared theorem surface.
Third, the fixed-cutoff measurement/Born-rule package is closed. Fourth, the
checkpoint/restoration observer package is closed. The external burdens on this lane are the final
closure of the patch-net embedding and compact heat-kernel lifts, global consensus analysis,
refinement-stable universality, and the continuum/gravity lift.

\paragraph{Reference-architecture takeaway.}
The synthesis-level point is that OPH has a concrete screen-side habitat in which one can
write local observables, overlap observables, record registers, and synchronization maps without
pretending to identify a unique final UV completion.

\paragraph{Theorem (Fixed-cutoff record algebra and Born-Luders package).}
For each completed compare/write/verify slice of the regulated microphysics, the declared pointer
and overlap-sector projectors generate a finite commutative record algebra
\[
\mathcal Z_{\mathrm{rec}}.
\]
Every observer-accessible event \(E\) in that algebra has probability
\[
\mathbb P(E)=\operatorname{Tr}(\rho P_E),
\]
and conditioning on \(E\) gives the operational post-measurement state
\[
\rho\!\mid_E = \frac{P_E \rho P_E}{\operatorname{Tr}(\rho P_E)}.
\]

\paragraph{Proof sketch.}
At fixed cutoff, the declared record and pointer projectors commute and the measurement surface
factors through that finite classical algebra. Finite-dimensional projective measurement on a
commuting projector algebra gives the Born trace and the Luders update directly.

\paragraph{Theorem (Checkpoint/restoration and observer backup).}
For an observer patch \(P_O\), let the checkpoint data consist of the observer-facing record
algebra, the accessible local state, the future update schedule, and the externally visible
overlap interface data. Exact restoration reproduces the full future law of observer-accessible
events, while an \(\varepsilon\)-accurate restoration changes that future law by at most
\(\varepsilon\) in total variation.

\paragraph{Proof sketch.}
Exact restoration reinstates the same accessible state on the same observer-facing algebra with
the same future channel sequence, so every later event probability agrees exactly. Approximate
restoration is controlled by contractivity of trace distance under completely positive
trace-preserving maps, which propagates the initial \(\varepsilon\) error bound to all later
observer-accessible events.

\paragraph{Observer-facing conclusion.}
The theorem-level content is modest but sharp: OPH supplies a fixed-cutoff measurement interface,
a fixed-cutoff checkpoint/restoration package, and an operational observer-identity criterion on
observer-accessible event algebras. Stronger substrate-selection or strange-loop closure claims
remain outside the recovered theorem package.

\paragraph{Extended derivation.}
The fixed-cutoff measurement, Born-rule, and checkpoint/restoration packages are developed in
\emph{Screen Microphysics, Patches, and Observer Synchronization in OPH}~\cite{ophmicrophysics}.

\section{Conditional Worldsheet/String Branch}

The synthesis paper carries the string/worldsheet branch at the same theorem boundary as the dedicated
string paper: a genuine theorem-level bridge from OPH edge sectors to the two-dimensional
Yang-Mills partition function, followed by a conditional large-\(N_{\mathrm{edge}}\) continuation
to a worldsheet reorganization.

\paragraph{Theorem (Heat-kernel edge-sector bridge).}
Under the OPH axioms, the fixed-cutoff overlap-gauge realization, and the local-Gibbs/MaxEnt edge
branch, the edge-sector weights satisfy
\[
p_R(t)\propto d_R e^{-t C_2(R)}.
\]
Consequently the edge partition function matches the standard two-dimensional Yang-Mills
heat-kernel form.

\paragraph{Proof sketch.}
The fixed-cutoff collar package produces a sector decomposition on the edge algebra, and the
local-Gibbs MaxEnt branch fixes the sector weights by the quadratic Casimir. Peter--Weyl
decomposition~\cite{peterWeyl27}
then identifies the resulting sum with the standard heat-kernel expression for
two-dimensional Yang-Mills.

\paragraph{Continuation boundary.}
If a controlled large-\(N_{\mathrm{edge}}\) regime exists, with \(N_{\mathrm{edge}}\) distinct from
the physical color rank \(N_c=3\), then the Gross-Taylor rewriting~\cite{grossTaylor93} gives a genus expansion and a
worldsheet-like reorganization. Critical-superstring claims, worldsheet CFT closure, anomaly
cancellation, and full massless-spectrum matching remain outside the recovered OPH core.

\paragraph{Extended derivation.}
This section records the synthesis-level statement; the longer technical discussion remains
in the repository string fragment.

\section{Cross-Lane Theorem Boundary and Open Directions}

The state of the OPH suite is:

\begin{enumerate}[leftmargin=2em]
\item The fixed-cutoff collar, higher-gauge, and consensus theorem packages are part of
the closed finite-regulator surface. The genuinely noncentral topological branch is not
fixed-cutoff gap, and the quotient-level repair normal form is explicit on the finite patch-net
surface.
\item The screen-microphysics lane has an explicit regulated reference architecture together
with closed measurement/Born-rule and checkpoint/restoration packages. Its remaining open fronts
are the final patch-net embedding pass, the compact heat-kernel lift, refinement-stable
universality, global consensus closure, and the continuum/gravity lift.
\item Physical UV uniqueness is quotient-level and ancilla-stable: OPH fixes the physical branch
modulo implementation hiding and inert ancillary refinement, not a unique microscopic
representative.
\item The Lorentz/null-modular/Einstein chain remains a controlled refinement/scaling-limit
branch. In the canonical premise ledger, the external burdens are the geometric modular
branch, the scaling-limit scope clause, and the fixed-cap stationarity step.
\item Compact gauge reconstruction and the realized Standard Model quotient remain tied to the
transportability, symmetric-braiding, and fiber-functor premises stated in the canonical
technical-premise list.
\item The particle branch is structurally rich and sector-split: structural carriers and the
realized Standard Model skeleton are fixed, electroweak closure is in hand, the Higgs/top stage is
quantitatively emitted on its secondary branch, and the charged-lepton, quark, neutrino, and
hadron closures sit on explicit continuation fronts.
\item The string branch is a continuation of the heat-kernel edge-sector theorem, not a promoted
recovered-core result.
\end{enumerate}

\paragraph{Companion papers.}
Detailed derivations appear in Refs.~\cite{ophcompact,ophconsensus,ophmicrophysics,ophparticles}.

\section{Conclusion}

Observer-Patch Holography is presented here as a single synthesis surface. The OPH basis, the
conditional gravity and gauge branches, the consensus formulation, the particle-spectrum branch,
the regulated screen-microphysics lane, the fixed-cutoff measurement and observer package, and the
conditional worldsheet continuation appear on one common synthesis surface.

The suite is uneven but concrete. OPH has a sharp
fixed-cutoff local picture, a conditional recovered core for relativity and Standard-Model
structure, real particle-side quantitative outputs, and a usable microscopic engineering lane. The
main external burdens are equally explicit: the continuum BW/geometric lift, the
transportable-sector compact-gauge lift, the remaining microphysics closure passes, and the open
matter-family and hadron closures. Refs.~\cite{ophcompact,ophconsensus,ophmicrophysics,ophparticles} remain the depth surfaces for
sector-by-sector proofs, calibration details, and continuation-level material.


\appendix
\input{tex_fragments/OBSERVERS_APPENDICES.tex}
\input{tex_fragments/OBSERVERS_REFERENCES.tex}

\end{document}
